\section{Related Work}
\label{sec:narrative-related}

\subsection{State Estimation and Physical Source Inference}
Unknown-input and DAE estimators recover electrical states despite
incomplete models or exogenous inputs
\cite{gillijns2007input,nugroho2022dae,katanic2024recursive}. Functional
observability asks whether a requested state quantity can be recovered
\cite{fernando2010functional,montanari2022functional}. Source inference asks
which intervention can explain the trajectory. Descriptor-model
diagnosability and nuisance-profiled tests formalize this distinction
\cite{eriksson2013diagnosability,scharf1994matched}.

Physical hypotheses support outage inference from phasor angles
\cite{tate2008outage}, Bayesian inference with uncertain states
\cite{zhao2014outages}, and sequential detection under stochastic operation
\cite{chen2016quickest,rovatsos2017outage}. Anguluri \emph{et al.} recover
sparse forcing locations, inputs and an unknown initial state using group
LASSO for linear dynamics \cite{anguluri2023group}. These are precedents
for source recovery; group LASSO is not an executed comparator in our
frozen campaign. Sensor placement also affects source distinguishability
\cite{jamei2020limits,conte2023fault,yildiz2025sparse}.

\subsection{Trajectory-Based and Sequential Diagnosis}
Warner \emph{et al.} compare simulated voltage/current fault trajectories
over a moving onset window and return candidate lines
\cite{warner2023highfidelity}. Their evaluated target is short-circuit
identification before clearing; load changes are an adverse-event test.
Our target follows load changes and line openings through restoration and
successive interventions. Physical trajectory libraries, remote PMUs and
retrospective onset scanning are shared ingredients. Their procedure was
not run on our data, so no numerical superiority is inferred.

The Adaptive Matrix CuSum derives sequential identification error/delay
results under specified observation laws \cite{warner2024matrix}.
Multiple-model filtering already propagates conditional state/mode
hypotheses \cite{blom1988imm}. The multiple-model multiple-hypothesis
filter maintains mode-sequence histories and adapts their number through
pruning and merging \cite{boers2005m3h}. Retaining histories or adapting
bank size is therefore not an original principle of our method.
Our pruned history bank uses these ideas,
but its approximate covariances, proposal screen and fitted startup context
do not inherit their general guarantees. The complete procedure evaluated
here makes retained recovery compete with later interventions and reports
timely attribution and alarm burden. These comparisons do not isolate
recovery memory from the other admission and publication rules.

\subsection{Learned Diagnosis and Physical Uncertainty}
Learned and knowledge-driven representations identify PMU disturbances and
assets \cite{pandey2020realtime,li2024knowledge,yildiz2024few}. We compare
two fixed ExtraTrees adapters from the SGSMA project. Their performance
tests transfer of those implementations, not the limits of all supervised
learning. A known physical catalog and source-excluded supervised training
also impose different information constraints.

Physical stochastic models support latent-state estimation
\cite{ensaf2023kernel,tartakovsky2023gp}, while joint DAE parameter inference
addresses uncertain generators and networks \cite{albustami2026dae}.
Our stream method uses a finite bank of constant network hypotheses.
Trajectory sensitivities and Gaussian evidence have also been combined
for Bayesian parameter inference in power-system DAEs
\cite{nagi2022trajectory}. Their use alone does not distinguish an event
attribution algorithm. The separate signature test examines a local event--context derivative;
it does not identify every plant parameter or replace the stream method
with an unrestricted joint estimator.

\begin{table*}[t]\centering\small
\caption{Closest inference contracts. Only Chen and the fixed ExtraTrees
adapters are executed comparators in the confirmatory stream campaign.
An untested capability is not an algorithmic impossibility.}
\label{tab:literature-capability}
\begin{tabularx}{\textwidth}{@{}lXXX@{}}\toprule
Route & Mechanism & Evaluated target & Role in this study\\\midrule
Chen / Rovatsos & Network likelihoods and sequential tests & Line outages; Rovatsos includes transients & Chen on persistent-line first decisions\\
Warner high-fidelity & Simulated V/I paths and onset search & Short-circuit line candidates & Closest trajectory precedent; no common-data run\\
Adaptive Matrix CuSum & Sequential diagnosis with error/delay analysis & One unknown distributional change & Statistical precedent; no inherited guarantee\\
DAE / stochastic-state estimation & Physical state and uncertainty models & States, unknown inputs or parameters & Components and adjacent tasks\\
Learned PMU diagnosis & Features or constrained representations & Event families and locations & Two frozen ExtraTrees references\\
Physical history bank here & Estimated-past-state paths and recovery & Successive catalog load/line events & Timely attribution, extras and restoration\\\bottomrule
\end{tabularx}\end{table*}
